\documentclass[11pt]{article}

\usepackage[margin=1in]{geometry}
\usepackage{amsmath,amssymb,amsthm}
\usepackage{bm}
\usepackage{hyperref}
\usepackage{enumitem}

\title{LST $\times$ Multiplicity:\\
Viability Kernels, Ethical Sub-Kernels, and Constraint Convergence}
\author{Research Note}
\date{\today}

\begin{document}
\maketitle

\begin{abstract}
This report distills a design conversation between Living Systems Tech (LST) and Multiplicity Theory into a precise kernel-based interface contract. The central tension is between internal mathematical invariance and external, heterogeneous constraint fields (biological, environmental, institutional, jurisdictional, lived). The outcome is a falsifiable mapping between LST's stress--capacity semantics and a prime-indexed multiplicity framework endowed with a viability kernel \(K\), an ethical sub-kernel \(K_{\text{ethical}}\subset K\), and a composite constraint operator \(F = C_{\text{inst}}\circ C_{\text{env}}\circ C_{\text{bio}}\circ T\). Bare dynamical stability (membership in \(K\)) is explicitly separated from ethically-endorsable stability (membership in \(K_{\text{ethical}}\)), with dissociative but functionally stable regimes characterized as \(K\setminus K_{\text{ethical}}\). The report formalizes these notions using per-atom Lyapunov functionals, small-gain inequalities, and proximal operators, and sketches how the resulting structure can be instantiated in a MultiplicityCell-style implementation. [file:5][web:71][web:72]
\end{abstract}
\newpage
\tableofcontents
\newpage
\section{Executive Summary}

\subsection{Context and Tension}

The discussion begins from a practical systems concern: most recursive validation frameworks preserve invariance under internally defined constraints but degrade when exposed to real-world, cross-domain stress fields. LST explicitly models stress versus capacity across biological, environmental, and institutional domains, treating \emph{truth} as whether a system holds under converging constraints rather than under idealized transformations. [file:5][web:77]

Multiplicity Theory and its kernel formalism approach the same problem from the opposite direction: they first define a mathematically curated ``lawful'' subspace and a contraction-based recursion, then ask which invariants survive across scales and perturbations. The central tension is therefore:
\begin{quote}
\emph{Can a prime-indexed, recursion-based invariance be made answerable to heterogeneous, non-commuting real-world constraint fields without collapsing back into a purely idealized model?}
\end{quote}
[file:5][web:74]

\subsection{Outcome: A Kernel-Based Interface Contract}

The conversation converges on a shared structure:

\begin{itemize}[leftmargin=*]
  \item A viability kernel \(K\subset X\) of multiplicity states that is:
  \begin{itemize}
    \item closed under the internal recursion \(T\),
    \item stable (on a specified time scale) under the composite constraint operator
    \[
      F = C_{\text{inst}}\circ C_{\text{env}}\circ C_{\text{bio}}\circ T,
    \]
    \item and nontrivially intersecting the set of non-collapsed lived configurations identified by LST. [file:5]
  \end{itemize}
  \item An ethical sub-kernel \(K_{\text{ethical}}\subset K\) defined by additional invariants corresponding to LST's notion of \emph{survives without self-destruction or dissociation}. [file:5]
  \item An operational notion of instability: for some \(x_0\in K\), iterates \(x_{t+1} = F(x_t)\) either leave \(K\) altogether (collapse) or leave \(K_{\text{ethical}}\) while remaining in \(K\) (ethically failed stability). [file:5]
\end{itemize}

Prime-indexed channels in the multiplicity state space are explicitly typed as biological load, environmental/infrastructural stress, institutional/jurisdictional constraint, and internal regulation/agency. Per-atom invariants split into bare viability and ethical components, and small-gain / contraction certificates (e.g., SlopeUB, GapLB) are interpreted as capacity constraints. Proximal operators implement both hard capacity bounds and ethical boundaries as generalized projections. [file:5][web:72][web:78]

\subsection{Ethical vs Bare Viability}

A key precision decision is made explicit:

\begin{itemize}[leftmargin=*]
  \item Deep, role-locked, institutionally rewarded dissociation is treated as:
  \[
    x_t \in K \setminus K_{\text{ethical}},
  \]
  i.e., dynamically viable but ethically failed.
  \item Collapse corresponds to trajectories leaving \(K\).
  \item Endorsed holding under pressure corresponds to trajectories remaining within \(K_{\text{ethical}}\). [file:5]
\end{itemize}

This separation ensures that the mathematics can distinguish:
\begin{enumerate}[label=(\roman*), leftmargin=*]
  \item systems that are dynamically stable but ethically unacceptable;
  \item systems that are both dynamically and ethically viable;
  \item systems that have genuinely failed in both senses.
\end{enumerate}

\section{Mathematical Overview}

\subsection{State Space and Channel Typing}

Let \(X\) denote the multiplicity state space. In the underlying theory, this can be thought of as a prime-indexed Hilbert space, but for the purposes of the kernel formalism we focus on the following structure:

\begin{itemize}[leftmargin=*]
  \item A set of \emph{atoms} indexed by primes \(p\in \mathbb{P}\).
  \item For each atom \(\pi\) (associated to some prime \(p\)), a coefficient vector \(c_\pi\in \mathbb{R}^{d_\pi}\).
  \item Each coordinate of \(c_\pi\) is typed as belonging to one of:
  \begin{enumerate}[label=(\alph*), leftmargin=*]
    \item biological load channels,
    \item environmental/infrastructural stress channels,
    \item institutional/jurisdictional constraint channels,
    \item internal regulation/agency channels.
  \end{enumerate}
  [file:5]
\end{itemize}

Thus a global state is \(x = (c_\pi)_{\pi\in \Pi}\in X\), where \(\Pi\) indexes the atoms.

\subsection{Internal Recursion}

The internal dynamics are given by a recursion \(T:X\to X\) defined atomwise via a proposal plus proximal step:

\begin{equation}
  c^{t+1}_\pi = \operatorname{Prox}_\pi\bigl(c^t_\pi + \operatorname{proposal}_\pi(c^t_\pi,\text{neighbors}); \text{budget}_\pi\bigr),
  \label{eq:atom-update}
\end{equation}

where:

\begin{itemize}[leftmargin=*]
  \item \(\operatorname{proposal}_\pi\) encodes the raw update from internal interactions and cross-atom couplings;
  \item \(\operatorname{Prox}_\pi(\cdot;\text{budget}_\pi)\) is a proximal operator that enforces both capacity and ethical constraints at atom \(\pi\); [file:5][web:72][web:75]
  \item \(\text{budget}_\pi\) is a parameter (or set of parameters) corresponding to the local capacity for that atom.
\end{itemize}

Global notation:
\[
  x^{t+1} = T(x^t),
\]
with \(T\) assembled from \eqref{eq:atom-update} across all atoms.

\subsection{Constraint Operators}

Real-world constraints are modeled as operators on \(X\):

\begin{align*}
  C_{\text{bio}} &: X \to X \quad \text{(biological stress)},\\
  C_{\text{env}} &: X \to X \quad \text{(environmental/infrastructural stress)},\\
  C_{\text{inst}} &: X \to X \quad \text{(institutional/jurisdictional constraints)}.
\end{align*}

These operators are generally non-commuting, reflecting LST's emphasis on path dependence (e.g., institutional trauma followed by environmental stress is not equivalent to the reverse ordering). They act primarily on their designated channel types but can have cross-channel effects, e.g., institutional stress affecting internal regulation channels. [file:5][web:74]

The combined dynamical map under constraint convergence is:

\begin{equation}
  F := C_{\text{inst}} \circ C_{\text{env}} \circ C_{\text{bio}} \circ T.
  \label{eq:F}
\end{equation}

Trajectories under real-world stress are given by:

\[
  x^{t+1} = F(x^t),\quad x^0\in X.
\]

\subsection{Viability Kernel \(K\)}

The kernel formalism defines a viability kernel \(K\subset X\) with three roles:

\begin{enumerate}[label=(\roman*), leftmargin=*]
  \item It is the domain on which the internal recursion \(T\) is contractive.
  \item It is invariant (or sufficiently stable) under \(F\) on the time scale of interest.
  \item It intersects nontrivially with the lived configurations LST identifies as non-collapsed systems. [file:5]
\end{enumerate}

Mathematically, this is enforced via Lyapunov and small-gain machinery. For each atom \(\pi\), we define a local Lyapunov functional \(V_\pi:X\to[0,\infty)\) and a global functional \(V:X\to[0,\infty)\), typically a weighted sum:

\[
  V(x) = \sum_{\pi} w_\pi V_\pi(x),
\]
for some weights \(w_\pi >0\). [file:5][web:71][web:77]

The key conditions are:

\begin{itemize}[leftmargin=*]
  \item \emph{Local contraction bound (SlopeUB):} for each atom \(\pi\), the local slope or gain of the update map satisfies an upper bound
  \[
    \text{SlopeUB}_\pi < 1,
  \]
  on states in \(K\). [file:5][web:80]
  \item \emph{Global Lyapunov decrease or boundedness:} for all \(x\in K\),
  \[
    V(F(x)) - V(x) \leq 0
  \]
  or more generally \(V(F(x)) \leq \alpha(V(x))\) for some \(\alpha\) with \(\alpha(s) < s\) away from the origin. [file:5][web:71][web:74]
  \item \emph{Non-collapse conditions:} additional constraints ensure that certain critical channels do not vanish or blow up, and that a minimum diversity of channels is preserved. [file:5]
\end{itemize}

Operationally, we can define \(K\) as the largest subset of \(X\) on which:

\begin{equation}
  \text{SlopeUB}_\pi(x) < 1,\quad \forall \pi,\quad V(x) \leq V_{\max},\quad \text{and non-collapse conditions hold.}
  \label{eq:K-def}
\end{equation}

Here \(V_{\max}\) can be chosen in accordance with LST's capacity thresholds.

\subsection{Ethical Sub-Kernel \(K_{\text{ethical}}\)}

Beyond bare viability, LST introduces an ethical dimension: some states are dynamically stable but ethically unacceptable (e.g., deep dissociation). The kernel note formalizes this through additional invariants and proximal constraints, yielding an ethical sub-kernel \(K_{\text{ethical}}\subset K\). [file:5]

Let \(E_\pi:X\to [0,\infty)\) denote per-atom ethical functionals, and define an aggregate ethical Lyapunov functional, e.g.,

\[
  E(x) = \sum_{\pi} \tilde{w}_\pi E_\pi(x).
\]

These functionals measure phenomena such as:

\begin{itemize}[leftmargin=*]
  \item excessive flattening of certain channels (e.g., suppression of affect),
  \item pathological asymmetries (e.g., institutional channels saturating while internal-regulation channels vanish),
  \item self-nullification patterns where a subsystem stabilizes by erasing its own signal. [file:5]
\end{itemize}

The proximal operators in \eqref{eq:atom-update} incorporate these ethical constraints by projecting onto sets that forbid such configurations. In particular, the ethical kernel is defined as:

\begin{equation}
  K_{\text{ethical}} := \left\{ x\in K \;\middle|\; E(x) \leq E_{\max},\ \text{and all ethical proximal constraints are satisfied} \right\}.
  \label{eq:K-ethical}
\end{equation}

The conversation explicitly agrees that:

\begin{itemize}[leftmargin=*]
  \item Deep, role-locked, institutionally rewarded dissociation lies in \(K\setminus K_{\text{ethical}}\) by default.
  \item Only states in \(K_{\text{ethical}}\) represent \emph{endorsed holding under pressure}.
  \item States in \(K\setminus K_{\text{ethical}}\) are dynamically viable but ethically failed, and states outside \(K\) are collapsed. [file:5]
\end{itemize}

\subsection{Instability as Trajectory Escape}

Instability is operationalized in terms of trajectory escape from these kernels. For a given initial state \(x_0\in K\), we consider the trajectory under \(F\):

\[
  x_{t+1} = F(x_t),\quad t=0,1,2,\dotsc
\]

Then:

\begin{itemize}[leftmargin=*]
  \item \emph{Bare viability} holds on a time horizon \(T\) if \(x_t\in K\) for all \(0\leq t\leq T\). [file:5]
  \item \emph{Ethical viability} holds on horizon \(T\) if \(x_t\in K_{\text{ethical}}\) for all \(0\leq t\leq T\). [file:5]
  \item \emph{Dissociative survival} occurs if there exists \(t\) with \(x_t\in K\setminus K_{\text{ethical}}\), even if the trajectory never leaves \(K\). [file:5]
  \item \emph{Collapse} occurs if there exists \(t\) with \(x_t\notin K\), i.e., if bare Lyapunov or non-collapse invariants fail. [file:5]
\end{itemize}

This structure allows direct alignment with LST's classification of outcomes:

\begin{center}
\begin{tabular}{l|l}
  LST classification & Kernel region \\
  \hline
  Endorsed holding under pressure & \(x_t\in K_{\text{ethical}}\ \forall t\) \\
  Survives by distortion / dissociation & \(x_t\in K\setminus K_{\text{ethical}}\) for some \(t\) \\
  Fails / collapses & \(x_t\notin K\) for some \(t\) \\
\end{tabular}
\end{center}
[file:5]

\section{From LST Variables to Kernel Inequalities}

\subsection{LST Stress and Capacity}

LST works with variables such as:

\begin{itemize}[leftmargin=*]
  \item biological load (trauma, fatigue, physiological stress),
  \item environmental stress (resource scarcity, climate, spatial constraints),
  \item institutional pressure (policy, surveillance, role conflict),
  \item internal regulation (self-regulation capacity, agency). [file:5][web:36][web:77]
\end{itemize}

Each variable typically comes with thresholds that divide regimes, e.g.:

\begin{itemize}[leftmargin=*]
  \item ``operationally sustainable'',
  \item ``functional but dissociated'',
  \item ``breakdown/collapse''.
\end{itemize}

The kernel mapping prescribes that each such threshold be translated into:

\begin{enumerate}[label=(\alph*), leftmargin=*]
  \item An inequality on channel magnitudes or derived quantities (e.g., norms, averages).
  \item Or an inequality on Lyapunov or ethical functionals (e.g., \(V(x)\leq V_{\max}\), \(E(x)\leq E_{\max}\)).
  \item Or a condition on SlopeUB/GapLB parameters representing small-gain bounds. [file:5][web:71]
\end{enumerate}

\subsection{Mapping Procedure}

For a given LST scenario, the workflow on the kernel side is:

\begin{enumerate}[label=\arabic*., leftmargin=*]
  \item \textbf{Extract variables and thresholds:} identify LST stress variables, capacity thresholds, and tipping points. [file:5]
  \item \textbf{Build multiplicity state space \(X\):} allocate prime-indexed channels for each stress class plus internal regulation. [file:5]
  \item \textbf{Define \(C_{\text{bio}}, C_{\text{env}}, C_{\text{inst}}\):} implement each LST stress source as an operator on designated channels, preserving non-commutation when appropriate. [file:5]
  \item \textbf{Implement \(T\):} use the kernel update \eqref{eq:atom-update} with per-atom proximal maps chosen so their invariants correspond to LST capacity notions (bounded load, no collapse, preservation of certain structures). [file:5][web:72]
  \item \textbf{Compute or specify \(K\):} define \(K\) via SlopeUB < 1, Lyapunov bounds, capacity inequalities, and non-collapse constraints. [file:5][web:71]
  \item \textbf{Define \(K_{\text{ethical}}\):} add ethical proximal constraints and Lyapunov terms (e.g., forbidding degenerate flats, enforcing balance) to obtain \(K_{\text{ethical}}\subset K\). [file:5]
  \item \textbf{Align LST thresholds with kernel regions:} show that LST's ``within capacity'', ``distorted survival'', and ``collapse'' correspond exactly to membership in \(K_{\text{ethical}}\), \(K\setminus K_{\text{ethical}}\), and \(X\setminus K\) respectively. [file:5]
\end{enumerate}

The result is a falsifiable mapping: if an LST-defined stress threshold is crossed while the kernel conditions still indicate membership in \(K_{\text{ethical}}\), the model is mis-specified and must be corrected.

\section{Proximal Maps and Small-Gain Conditions}

\subsection{Proximal Operators as Capacity and Ethics}

Proximal operators play a central role in implementing both capacity and ethical constraints in the kernel update. In convex analysis, the proximal operator of a proper convex function \(f:\mathbb{R}^n\to \mathbb{R}\cup\{+\infty\}\) with parameter \(t>0\) is defined as:

\[
  \text{prox}_{t f}(v) = \arg\min_{x\in\mathbb{R}^n} \left( f(x) + \frac{1}{2t}\|x - v\|_2^2 \right).
\]

This generalizes projection onto a convex set and can be interpreted as a generalized gradient step that stays close to \(v\). [web:72][web:78]

In the kernel, \(\operatorname{Prox}_\pi(\cdot;\text{budget}_\pi)\) is designed so that:

\begin{itemize}[leftmargin=*]
  \item capacity constraints are encoded as indicator functions of feasible sets (e.g., norm bounds, inter-channel inequalities), so their proximal maps are projections onto those sets; [file:5][web:75]
  \item ethical constraints are encoded via additional penalty terms and indicator functions enforcing non-dissociation conditions (e.g., lower bounds, balance constraints), so that their proximal maps disallow self-nullifying solutions. [file:5]
\end{itemize}

Thus the update \eqref{eq:atom-update} can be seen as computing a local proposal, then projecting it back into the intersection of capacity and ethical feasibility regions at each atom.

\subsection{Small-Gain and Lyapunov Conditions}

Small-gain theorems and Lyapunov-based analysis provide conditions under which interconnected subsystems remain stable. In hybrid and nonlinear settings, such theorems use Lyapunov functions and gain conditions to guarantee that the feedback interconnection is globally asymptotically stable. [web:71][web:74][web:80]

In the kernel, these ideas are specialized as:

\begin{itemize}[leftmargin=*]
  \item \emph{Per-atom small-gain bounds} (SlopeUB and GapLB) ensuring that the local effect of updates and couplings is contractive on average.
  \item \emph{Global Lyapunov conditions} ensuring that the sum of local Lyapunov functionals decreases or remains bounded under \(F\).
  \item \emph{Ethical Lyapunov} terms ensuring that ethically undesirable configurations incur non-decreasing cost unless corrected by proximal maps. [file:5]
\end{itemize}

Combined, these conditions provide a mathematically grounded version of LST's ``stress not exceeding capacity'' heuristic: exceeding capacity translates into violating small-gain or Lyapunov inequalities, causing escape from \(K\) or \(K_{\text{ethical}}\).

\section{Illustrative Code Skeleton}

To connect the formalism to implementation, we sketch a high-level pseudocode of a MultiplicityCell-like update that respects the kernel structure. This is schematic and omits many details, but shows where \(K\) and \(K_{\text{ethical}}\) live in code. [file:5]

\subsection{Typed Atom Update}

\begin{verbatim}
// Pseudocode / high-level sketch

struct AtomState {
  Vector bio;      // biological load channels
  Vector env;      // environmental stress channels
  Vector inst;     // institutional constraint channels
  Vector internal; // internal regulation/agency channels
}

struct KernelParams {
  CapacityBounds capacity;   // per-channel capacity constraints
  EthicalBounds ethical;     // ethical invariants (non-dissociation, balance)
  LyapunovWeights V_weights; // weights for bare Lyapunov V
  LyapunovWeights E_weights; // weights for ethical Lyapunov E
}

AtomState proposal(AtomState c_pi, NeighborStates nb) {
  // Compute raw proposal from internal dynamics and couplings
  // (details depend on the specific multiplicity model)
}

AtomState prox_capacity(AtomState v, CapacityBounds cap) {
  // Project v onto capacity-feasible region:
  //  - norm bounds on channels
  //  - inter-channel inequalities representing "within capacity"
}

AtomState prox_ethics(AtomState v, EthicalBounds eth) {
  // Project v onto ethically-feasible region:
  //  - forbid degenerate flats (dissociation patterns)
  //  - enforce balance between selected channels
}

AtomState update_atom(AtomState c_pi, NeighborStates nb, KernelParams params) {
  AtomState prop = proposal(c_pi, nb);
  AtomState v = c_pi + prop; // raw step
  
  // Capacity projection (bare viability)
  AtomState v_cap = prox_capacity(v, params.capacity);
  
  // Ethical projection (ethical viability)
  AtomState v_eth = prox_ethics(v_cap, params.ethical);
  
  return v_eth;
}
\end{verbatim}

In this pattern:

\begin{itemize}[leftmargin=*]
  \item \texttt{prox\_capacity} enforces membership in \(K\).
  \item \texttt{prox\_ethics} enforces membership in \(K_{\text{ethical}}\).
\end{itemize}

If one wishes to distinguish trajectories that are merely in \(K\) from those in \(K_{\text{ethical}}\), one can track pre-ethical and post-ethical states separately.

\subsection{Global Step and Classification}

\begin{verbatim}
struct TrajectoryStatus {
  bool in_K;           // bare viability
  bool in_K_ethical;   // ethical viability
}

TrajectoryStatus classify_state(State x, KernelParams params) {
  double V = compute_bare_lyapunov(x, params.V_weights);
  double E = compute_eth_lyapunov(x, params.E_weights);
  
  bool in_K = (V <= params.capacity.V_max) &&
              non_collapse_conditions(x);
  bool in_K_ethical = in_K &&
                      (E <= params.ethical.E_max) &&
                      ethical_invariants_satisfied(x);
  
  return { in_K, in_K_ethical };
}

void step(State& x, KernelParams params) {
  // Apply internal recursion T
  State x_T = apply_T(x, params);
  
  // Apply constraint operators
  State x_bio  = C_bio(x_T);
  State x_env  = C_env(x_bio);
  State x_inst = C_inst(x_env);
  
  x = x_inst; // This is F(x)
}
\end{verbatim}

A simulation loop can then implement:

\begin{verbatim}
State x = x0;
for (int t = 0; t < T; ++t) {
  auto status = classify_state(x, params);
  if (!status.in_K) {
    // collapse detected
    break;
  }
  if (!status.in_K_ethical) {
    // dissociative survival (ethically failed stability)
    // mark or intervene as needed
  }
  step(x, params);
}
\end{verbatim}

This illustrates how the conceptual distinctions among \(K\), \(K_{\text{ethical}}\), and their complements become concrete and testable in code. [file:5]

\section{Recommended Next Artifact: Scenario-001}

To fully instantiate the interface, the next artifact should be a joint document, \texttt{LSTxMultiplicity-Scenario-001.md}, containing:

\begin{itemize}[leftmargin=*]
  \item A narrative description (1--3 paragraphs) of a real LST scenario (level, key constraints, stress episodes).
  \item A list of LST variables and stress fields (names, units or qualitative scales, capacity thresholds).
  \item A short description of what counts as:
  \begin{itemize}
    \item endorsed holding under pressure,
    \item distorted survival / dissociation,
    \item collapse.
  \end{itemize}
  \item A mapping table from LST variables and thresholds to:
  \begin{itemize}
    \item multiplicity channels,
    \item kernel inequalities,
    \item classification regions \(K_{\text{ethical}}\), \(K\setminus K_{\text{ethical}}\), \(X\setminus K\).
  \end{itemize}
\end{itemize}

On the kernel side, this will enable a complete worked example in which LST stress thresholds and kernel membership conditions coincide exactly, turning the high-level design into a falsifiable, scenario-specific test. [file:5][web:73][web:70]

\section{Conclusion}

This report has turned a conceptual dialogue between LST and Multiplicity Theory into a concrete interface contract anchored in a viability kernel \(K\), an ethical sub-kernel \(K_{\text{ethical}}\), and a composite constraint operator \(F\). The key advance is the explicit separation between bare dynamical stability and ethically endorsed stability, realized via per-atom Lyapunov functionals, small-gain conditions, and proximal maps enforcing both capacity and ethical constraints. The next step is to instantiate this structure on a specific LST scenario and verify that the classification of trajectories by LST matches the kernel-based classification by membership in \(K_{\text{ethical}}\), \(K\setminus K_{\text{ethical}}\), or outside \(K\).

\appendix
\section{Mathematical Appendix}
\label{app:math}

This appendix collects the core proofs and operator norm bounds underpinning the viability kernel \(K\), the ethical sub-kernel \(K_{\text{ethical}}\), and the composite operator
\[
  F := C_{\text{inst}} \circ C_{\text{env}} \circ C_{\text{bio}} \circ T.
\]
The arguments are phrased in standard Lyapunov/small-gain and proximal calculus terms and adapted to the kernel structure developed in the main text. [file:5][web:71][web:72]

\subsection{Setup and Notation}

Let \(X\) be a Banach space with norm \(\|\cdot\|\), representing the global multiplicity state (assembled from prime-indexed atoms). For each atom \(\pi\) we have a local state \(c_\pi\) and a local norm \(\|\cdot\|_\pi\), with
\[
  \|x\|^2 = \sum_{\pi} \alpha_\pi \|c_\pi\|_\pi^2,
\]
for some weights \(\alpha_\pi>0\). [file:5]

The internal recursion \(T:X\to X\) is defined atomwise via
\begin{equation}
  c^{+}_\pi = \operatorname{Prox}_\pi\bigl(c_\pi + \operatorname{proposal}_\pi(c_\pi,\text{neighbors}); \text{budget}_\pi\bigr),
  \label{eq:app-prox-update}
\end{equation}
with \(x^+ = T(x)\) the global post-update state. The constraint operators
\[
  C_{\text{bio}},\ C_{\text{env}},\ C_{\text{inst}} : X \to X
\]
encode biological, environmental, and institutional stress respectively and are generally non-commuting. [file:5]

We use the usual operator norm
\[
  \|A\| := \sup_{\|x\|\neq 0} \frac{\|A x\|}{\|x\|},
\]
for any bounded linear operator \(A:X\to X\). [web:41]

\subsection{Nonexpansiveness of Proximal Steps}

We first recall a standard fact from convex analysis: proximal operators of proper, closed, convex functions are firmly nonexpansive. This property is central for viewing the per-atom proximal maps as contraction-like in the kernel. [web:72][web:97]

\begin{proposition}[Firm nonexpansiveness of proximal operators]
Let \(f:X\to\mathbb{R}\cup\{+\infty\}\) be proper, closed, and convex, and let
\[
  \operatorname{prox}_{t f}(v) := \arg\min_{x\in X} \left\{ f(x) + \frac{1}{2t}\|x-v\|^2 \right\},\quad t>0.
\]
Then for all \(x,y\in X\),
\[
  \|\operatorname{prox}_{t f}(x) - \operatorname{prox}_{t f}(y)\|^2 \leq
  \langle \operatorname{prox}_{t f}(x) - \operatorname{prox}_{t f}(y), x - y \rangle.
\]
In particular, \(\operatorname{prox}_{t f}\) is nonexpansive:
\[
  \|\operatorname{prox}_{t f}(x) - \operatorname{prox}_{t f}(y)\| \leq \|x - y\|.
\]
\label{prop:prox-nonexp}
\end{proposition}

\begin{proof}
This is a standard result; we sketch the argument. By optimality conditions, if \(u = \operatorname{prox}_{t f}(x)\) and \(v = \operatorname{prox}_{t f}(y)\), then
\[
  0 \in \partial f(u) + \frac{1}{t}(u - x), \qquad
  0 \in \partial f(v) + \frac{1}{t}(v - y).
\]
Thus
\[
  x - u \in t\,\partial f(u), \qquad y - v \in t\,\partial f(v).
\]
Monotonicity of the subdifferential \(\partial f\) gives
\[
  \langle (x - u) - (y - v), u - v \rangle \geq 0,
\]
which rearranges to
\[
  \langle u - v, x - y \rangle \geq \|u - v\|^2.
\]
This is the firm nonexpansiveness inequality; the nonexpansiveness bound follows by Cauchy-Schwarz. [web:72][web:97]
\end{proof}

\begin{corollary}[Nonexpansive per-atom proximal maps]
Suppose each \(\operatorname{Prox}_\pi\) in \eqref{eq:app-prox-update} is the proximal operator of a proper closed convex function with respect to \(\|\cdot\|_\pi\). Then for each \(\pi\),
\[
  \|\operatorname{Prox}_\pi(u) - \operatorname{Prox}_\pi(v)\|_\pi \leq \|u-v\|_\pi,\quad \forall u,v.
\]
In particular, the global map assembling all atoms is nonexpansive with respect to \(\|\cdot\|\). [file:5][web:72]
\end{corollary}

\subsection{Operator Norm Bound for the Composite Map \(F\)}

We now give an explicit operator norm bound for the composite map
\[
  F = C_{\text{inst}}\circ C_{\text{env}}\circ C_{\text{bio}}\circ T,
\]
under assumptions separating the local proposal dynamics from the proximal maps. [file:5]

\begin{assumption}[Local Lipschitz proposal]
For each atom \(\pi\), the proposal map \(P_\pi(x) := \operatorname{proposal}_\pi(c_\pi,\text{neighbors})\) is Lipschitz continuous in the global state \(x\), with constant \(L_\pi\), i.e.,
\[
  \|P_\pi(x) - P_\pi(y)\|_\pi \leq L_\pi \|x-y\|,\quad \forall x,y\in X.
\]
Let \(L_T := \sup_\pi L_\pi\). [file:5]
\end{assumption}

\begin{assumption}[Lipschitz constraint operators]
The constraint operators are Lipschitz with constants \(L_{\text{bio}}, L_{\text{env}}, L_{\text{inst}}\), i.e.,
\[
  \|C_{\text{bio}}(x) - C_{\text{bio}}(y)\| \leq L_{\text{bio}} \|x-y\|,
\]
and similarly for \(C_{\text{env}}, C_{\text{inst}}\). [file:5]
\end{assumption}

\begin{lemma}[Lipschitz bound for \(T\)]
Under the previous assumptions and the nonexpansiveness of per-atom proximal maps, the global recursion \(T\) is Lipschitz with constant
\[
  L_T^{\text{eff}} \leq 1 + L_T.
\]
\label{lem:Lipschitz-T}
\end{lemma}

\begin{proof}
Consider two global states \(x,y\in X\). At atom \(\pi\),
\[
  c_\pi^+ = \operatorname{Prox}_\pi(c_\pi + P_\pi(x)),\quad
  d_\pi^+ = \operatorname{Prox}_\pi(d_\pi + P_\pi(y)),
\]
where \(c_\pi\) and \(d_\pi\) are the components of \(x\) and \(y\). Then
\begin{align*}
  \|c_\pi^+ - d_\pi^+\|_\pi
  &\leq \|(c_\pi + P_\pi(x)) - (d_\pi + P_\pi(y))\|_\pi
  &&\text{(nonexpansiveness of }\operatorname{Prox}_\pi\text{)}\\
  &= \| (c_\pi - d_\pi) + (P_\pi(x) - P_\pi(y)) \|_\pi \\
  &\leq \|c_\pi - d_\pi\|_\pi + \|P_\pi(x) - P_\pi(y)\|_\pi \\
  &\leq \|c_\pi - d_\pi\|_\pi + L_\pi \|x-y\|.
\end{align*}
Summing over atoms with weights \(\alpha_\pi\) and using the definition of the global norm,
\[
  \|T(x) - T(y)\|
  \leq \|x-y\| + L_T \|x-y\|
  = (1 + L_T)\|x-y\|.
\]
Thus \(T\) is Lipschitz with constant at most \(1+L_T\). [file:5]
\end{proof}

\begin{theorem}[Global contraction bound for \(F\)]
Suppose:
\begin{enumerate}[label=(\roman*), leftmargin=*]
  \item The per-atom proximal maps are nonexpansive.
  \item The proposal Lipschitz constant satisfies \(L_T < \infty\).
  \item The constraint operators satisfy
  \[
    \|C_{\text{bio}}\|\leq L_{\text{bio}},\quad
    \|C_{\text{env}}\|\leq L_{\text{env}},\quad
    \|C_{\text{inst}}\|\leq L_{\text{inst}}.
  \]
\end{enumerate}
Then the composite map \(F\) is Lipschitz with constant
\[
  L_F \leq L_{\text{inst}}\,L_{\text{env}}\,L_{\text{bio}}\,(1 + L_T).
\]
In particular, if
\[
  L_{\text{inst}}\,L_{\text{env}}\,L_{\text{bio}}\,(1 + L_T) < 1,
\]
then \(F\) is a contraction on \(X\). [file:5][web:41][web:94]
\label{thm:F-contraction}
\end{theorem}

\begin{proof}
We have
\begin{align*}
  \|F(x) - F(y)\|
  &= \|C_{\text{inst}}(C_{\text{env}}(C_{\text{bio}}(T(x))) - C_{\text{env}}(C_{\text{bio}}(T(y))) )\| \\
  &\leq L_{\text{inst}} \|C_{\text{env}}(C_{\text{bio}}(T(x))) - C_{\text{env}}(C_{\text{bio}}(T(y)))\| \\
  &\leq L_{\text{inst}} L_{\text{env}} \|C_{\text{bio}}(T(x)) - C_{\text{bio}}(T(y))\| \\
  &\leq L_{\text{inst}} L_{\text{env}} L_{\text{bio}} \|T(x) - T(y)\|.
\end{align*}
Using Lemma~\ref{lem:Lipschitz-T},
\[
  \|T(x) - T(y)\| \leq (1+L_T)\|x-y\|,
\]
hence
\[
  \|F(x) - F(y)\| \leq
  L_{\text{inst}} L_{\text{env}} L_{\text{bio}} (1+L_T)\|x-y\|.
\]
The stated Lipschitz bound follows. If the product \(L_{\text{inst}} L_{\text{env}} L_{\text{bio}} (1+L_T) < 1\), then \(F\) is a contraction. [file:5][web:41][web:96]
\end{proof}

\begin{corollary}[Existence and uniqueness of a fixed point]
If \(X\) is complete and \(L_F<1\), then \(F\) has a unique fixed point \(x^\star\in X\), and for any initial state \(x^0\in X\),
\[
  \|F^t(x^0) - x^\star\| \leq L_F^t \|x^0 - x^\star\|\to 0.
\]
\label{cor:Banach-F}
\end{corollary}

\begin{proof}
This is the Banach fixed-point theorem applied to \(F\). [web:41][web:102]
\end{proof}

\subsection{Lyapunov-Based Kernel Characterization}

We now connect the operator norm bounds to the Lyapunov characterization of \(K\) and \(K_{\text{ethical}}\) used in the main text. [file:5][web:71]

\begin{assumption}[Per-atom Lyapunov functions]
For each atom \(\pi\), there exists a continuous function \(V_\pi:X\to[0,\infty)\) such that:
\begin{enumerate}[label=(\alph*), leftmargin=*]
  \item \(V_\pi(x) = 0\) if and only if the local state for \(\pi\) is at its desired equilibrium (or in a designated neutral set).
  \item There exist class-\(\mathcal{K}_\infty\) functions\footnote{Standard comparison functions as in ISS and small-gain theory.} \(\alpha_{1,\pi}, \alpha_{2,\pi}\) satisfying
  \[
    \alpha_{1,\pi}(\|c_\pi\|_\pi) \leq V_\pi(x) \leq \alpha_{2,\pi}(\|c_\pi\|_\pi).
  \]
\end{enumerate}
\end{assumption}

Define the global Lyapunov function
\[
  V(x) := \sum_{\pi} w_\pi V_\pi(x),
\]
for weights \(w_\pi>0\). [file:5]

\begin{assumption}[Lyapunov decrease on \(K\)]
There exists a set \(K\subset X\) and a class-\(\mathcal{K}_\infty\) function \(\alpha_V\) such that for all \(x\in K\),
\[
  V(F(x)) - V(x) \leq -\alpha_V(\|x\|).
\]
\end{assumption}

\begin{theorem}[Invariance and attractivity of \(K\)]
Suppose:
\begin{enumerate}[label=(\roman*), leftmargin=*]
  \item \(K\subset X\) is forward invariant under \(F\), i.e., \(x\in K \Rightarrow F(x)\in K\).
  \item The Lyapunov decrease condition holds on \(K\).
  \item The level sets \(\{x: V(x)\leq c\}\) are compact for all \(c\geq 0\) (radial unboundedness).
\end{enumerate}
Then:
\begin{enumerate}[label=(a), leftmargin=*]
  \item Every trajectory starting in \(K\) remains in \(K\) and satisfies \(V(x^t)\) strictly decreasing until it reaches the largest invariant subset contained in \(\{x\in K: \alpha_V(\|x\|)=0\}\).
  \item If \(\alpha_V(\|x\|)=0\) only at a set of desired equilibria \(E\subset K\), then every trajectory starting in \(K\) converges to \(E\).
\end{enumerate}
\label{thm:K-Lyapunov}
\end{theorem}

\begin{proof}
This is a standard Lyapunov/LaSalle argument. The Lyapunov decrease ensures that on \(K\), \(V(x^t)\) is strictly decreasing except possibly where \(\alpha_V(\|x\|)=0\), hence bounded and convergent. Radial unboundedness ensures trajectories remain bounded. The invariance of \(K\) and LaSalle's invariance principle then imply convergence to the largest invariant set contained in \(\{x\in K: \alpha_V(\|x\|)=0\}\), which in the second case is the set \(E\) of desired equilibria. [web:71][web:81]
\end{proof}

\subsection{Ethical Sub-Kernel \(K_{\text{ethical}}\) via Additional Lyapunov Terms}

To capture ethical constraints, we introduce additional per-atom functionals \(E_\pi:X\to[0,\infty)\) measuring ethical distortion (e.g., dissociation, channel flattening), and define
\[
  E(x) := \sum_{\pi} \tilde{w}_\pi E_\pi(x),
\]
with \(\tilde{w}_\pi>0\). [file:5]

\begin{assumption}[Ethical monotonicity on \(K_{\text{ethical}}\)]
There exists a set \(K_{\text{ethical}}\subset K\) and a nonnegative function \(\beta_E\) such that for all \(x\in K_{\text{ethical}}\),
\[
  E(F(x)) - E(x) \leq -\beta_E(\|x\|),
\]
with \(\beta_E(s) > 0\) whenever \(E(x)\) exceeds a designated ethical threshold. [file:5]
\end{assumption}

\begin{definition}[Ethical sub-kernel]
Define the ethical kernel as
\[
  K_{\text{ethical}} := \{ x\in K : E(x) \leq E_{\max},\ \text{all ethical proximal constraints hold} \},
\]
where \(E_{\max}\) encodes the maximal ethically admissible distortion consistent with LST's ``endorsed holding'' regime. [file:5]
\end{definition}

\begin{theorem}[Classification via \((V,E)\)]
Let \(x^0\in K\) and \(x^{t+1} = F(x^t)\). Under the assumptions above:
\begin{enumerate}[label=(a), leftmargin=*]
  \item If \(x^t\in K_{\text{ethical}}\) for all \(t\) and \(E\) is non-increasing, the trajectory is both dynamically and ethically viable (endorsed holding).
  \item If \(x^t\in K\) for all \(t\) but there exists \(t^\ast\) with \(x^{t^\ast}\notin K_{\text{ethical}}\), the trajectory is dynamically stable but ethically failed (dissociative survival).
  \item If there exists \(t^\ast\) with \(x^{t^\ast}\notin K\), the trajectory is dynamically non-viable (collapse).
\end{enumerate}
\label{thm:classification}
\end{theorem}

\begin{proof}
(a) If \(x^t\in K_{\text{ethical}}\) for all \(t\), then by definition \(x^t\in K\) and both the bare Lyapunov and ethical constraints are maintained. The monotonicity of \(E\) ensures ethical distortion does not increase; combined with the Lyapunov decrease for \(V\), this yields convergence to a set of states that are both dynamically and ethically acceptable. [file:5][web:71]

(b) If \(x^t\in K\) but \(x^{t^\ast}\notin K_{\text{ethical}}\) for some \(t^\ast\), then by construction \(x^{t^\ast}\) violates at least one ethical constraint (either \(E(x^{t^\ast})>E_{\max}\) or a proximal ethical constraint fails). Thus the trajectory is still dynamically stable but lies in \(K\setminus K_{\text{ethical}}\), matching the definition of dissociative survival. [file:5]

(c) If \(x^{t^\ast}\notin K\) for some \(t^\ast\), then bare Lyapunov or non-collapse invariants have failed: either \(V(x^{t^\ast})>V_{\max}\) or a non-collapse condition is broken. This corresponds to LST's ``collapse'' regime. [file:5]
\end{proof}

\subsection{Operator Norm vs Lyapunov Conditions}

Finally, we relate the norm-based contraction condition of Theorem~\ref{thm:F-contraction} to the Lyapunov-based characterization in Theorem~\ref{thm:K-Lyapunov}. [file:5][web:94]

\begin{proposition}[From contraction to Lyapunov decrease]
Suppose \(F:X\to X\) is a contraction with constant \(L_F<1\). Define
\[
  V(x) := \frac{1}{2}\|x - x^\star\|^2,
\]
where \(x^\star\) is the unique fixed point of \(F\). Then
\[
  V(F(x)) - V(x) \leq -\frac{1-L_F^2}{2} \|x - x^\star\|^2
\]
for all \(x\in X\), i.e., the Lyapunov decrease condition of Theorem~\ref{thm:K-Lyapunov} holds globally with \(\alpha_V(s) = \frac{1-L_F^2}{2} s^2\). [web:41][web:99]
\end{proposition}

\begin{proof}
We compute
\begin{align*}
  V(F(x)) - V(x)
  &= \frac{1}{2}\|F(x) - x^\star\|^2 - \frac{1}{2}\|x - x^\star\|^2 \\
  &\leq \frac{1}{2} L_F^2 \|x - x^\star\|^2 - \frac{1}{2}\|x - x^\star\|^2 \\
  &= -\frac{1-L_F^2}{2} \|x - x^\star\|^2,
\end{align*}
where we used the contraction property for \(F\) and the fact that \(F(x^\star) = x^\star\). [web:41][web:99]
\end{proof}

\begin{corollary}[Embedding norm-based bounds into kernel conditions]
If the contraction bound of Theorem~\ref{thm:F-contraction} holds on a subset \(K\subset X\), then choosing \(V(x) = \frac{1}{2}\|x - x^\star\|^2\) and restricting to \(K\) yields a Lyapunov decrease condition compatible with Theorem~\ref{thm:K-Lyapunov}. Thus the norm-based kernel definition (via \(L_F<1\)) and the Lyapunov-based kernel definition (via \(V\)) can be aligned. [file:5][web:71]
\end{corollary}

\subsection{Summary}

The arguments above justify three key structural claims used in the main report:

\begin{itemize}[leftmargin=*]
  \item Nonexpansive proximal steps and bounded proposal dynamics yield explicit Lipschitz bounds for the internal recursion \(T\) and the composite map \(F\).
  \item Under suitable parameter choices, these bounds imply contraction on a kernel \(K\), guaranteeing existence and uniqueness of a fixed point and exponential convergence within \(K\).
  \item Augmenting the bare Lyapunov functional with ethical functionals produces an ethical sub-kernel \(K_{\text{ethical}}\subset K\), enabling precise classification of trajectories into endorsed holding, dissociative survival, and collapse, in alignment with LST semantics.
\end{itemize}

These results give the mathematical backbone for treating LST stress/capacity thresholds as kernel membership conditions in the multiplicity framework. [file:5][web:71][web:72]

\nocite{*}
\bibliographystyle{plainnat}
\bibliography{references}
\end{document}

